\clearpage
\thispagestyle{plain}

\begin{center}
    {\Large \bfseries FICHA RESUMEN}
\end{center}

\vspace{0.3cm}

\begin{tabular}{@{}ll}
\textbf{Grado:} & Licenciatura en Ingeniería de Sistemas \\
\textbf{Modalidad:} & Proyecto de Grado \\
\textbf{Facultad:} & Ciencias y Tecnología \\
\textbf{Carrera:} & Ingeniería de Sistemas \\
\textbf{Título del trabajo:} & \parbox[t]{10cm}{Sistema Inteligente de Entrenamiento en
Gimnasios con Realidad Aumentada e Inteligencia Artificial} \\
\textbf{Autora:} & \autorTesis \\
\textbf{Tutor:} & \tutorTesis \\
\end{tabular}

\vspace{0.4cm}

En el primer capítulo se realizó una investigación sobre los antecedentes y se
analizó la problemática relacionada con el abandono temprano de los usuarios de
gimnasios, atribuida principalmente a la falta de orientación, motivación y
seguimiento personalizado. A partir de ello, se propone el diseño de un sistema
inteligente capaz de guiar al usuario en uso correcto de máquinas, personalizar
rutinas de entrenamiento y registrar su progreso físico.

En el segundo capítulo se desarrolla el marco teórico que sustenta la solución
tecnológica. Se revisan distintas metodologías de desarrollo de software,
comparando enfoques tradicionales y ágiles, y se justifica la selección de
Mobile-D como la metodología más adecuada para este proyecto.

En el tercer capítulo se detalla el área de aplicación y se describe la
adopción de la metodología Mobile-D en el desarrollo del sistema, especificando
su aplicación práctica en fases de exploración, inicialización, producción,
estabilización y pruebas.
